
Large language models have demonstrated code generation capabilities on benchmarks such as HumanEval and APPS, yet they lack formal correctness guarantees. Test suites detect obvious errors but miss edge cases that static analysis and satisfiability modulo theories (SMT) solvers can identify. However, batch SMT verification does not scale well to iterative debugging workflows: re-verifying entire programs at each repair step becomes computationally prohibitive.

Prior SMT-guided program repair systems such as Angelix and Prophet target small human-written patches, not full LLM-generated programs. Neural code generation systems such as AlphaCode and CodeT5 rely exclusively on test-suite validation, providing no formal verification. This gap motivates the hypothesis that incremental SMT verification---re-verifying only modified functions and their dependencies---could reduce verification time for LLM repair workflows in statically-typed languages.

The hypothesis builds on established components: static analyzers (Pyre for Python, Prusti for Rust) extract SMT constraints from type annotations in human-written code, Z3 supports incremental solving via push/pop contexts, and neural repair literature indicates that 80\% of human code repairs touch 1--3 lines. However, transferability to LLM-generated code remained unverified.

This work designed a four-hypothesis validation chain: h-e1 (existence hypothesis: static analyzers can extract SMT constraints from LLM-generated typed code at $\geq 90$\% success rate), followed by three mechanism hypotheses testing constraint extraction (h-m1), dependency analysis (h-m2), and speedup measurement (h-m3). The experimental protocol required extending HumanEval with Pydantic type annotations to create a typed benchmark.

The validation attempt encountered an invalid Anthropic API key, blocking all 48 LLM generation attempts with HTTP 401 errors. Without generated code, the constraint extraction pipeline could not test whether type annotations enable SMT predicate mapping. The hypothesis remains theoretically plausible but empirically unvalidated.

Despite generation failure, the dataset extension pipeline succeeded: 100 Pydantic-annotated prompts were mechanically generated from HumanEval. The constraint extraction architecture (DatasetExtender $\rightarrow$ LLMGenerator $\rightarrow$ PyreExtractor $\rightarrow$ Z3Validator $\rightarrow$ MetricsEngine) was validated on error files, confirming that the pipeline processes malformed inputs without failure.

This paper contributes:

\begin{enumerate}
\item \textbf{HumanEval + Pydantic Type Extensions}: A mechanically-generated typed benchmark dataset of 100 prompts suitable for SMT constraint extraction research, addressing the gap in typed benchmarks for formal verification of neural code generation.
\end{enumerate}

\begin{enumerate}
\item \textbf{Constraint Extraction Pipeline Architecture}: A modular five-component pipeline validated on negative cases and ready for retry with valid API credentials.
\end{enumerate}

The primary methodological contribution is the demonstration that infrastructure reliability is a prerequisite for hypothesis validation. Experimental pipelines must address failure modes explicitly through pre-flight validation, fast-fail logic for non-retryable errors, and checkpoint/resume mechanisms, or risk perpetual blockage.

Section 2 positions this work against SMT-guided repair for human code and neural code generation. Section 3 describes the hypothesis decomposition, experimental design, and pipeline architecture. Section 4 presents results: dataset extension succeeded, LLM generation failed, hypothesis untested. Section 5 discusses implications and limitations. Section 6 concludes with recommendations for retry and future work.
